\documentclass[10pt,a4paper]{amsart}
\usepackage[margin=19mm]{geometry}
\usepackage{amsmath,amssymb,mathtools}
\usepackage[scr=rsfs,cal=euler]{mathalfa}
\usepackage{xfrac}
\usepackage[unicode,hidelinks,bookmarks=false]{hyperref}
\newcommand{\ie}{i.e.\ }
\newcommand{\cf}{cf.\ }
\newcommand{\resp}{respectively\ }
\newcommand{\Subsection}{\subsection}
\providecommand{\nobreakdash}{\nobreak\hbox{-}}
\setlength{\emergencystretch}{2em}
\multlinegap0pt
\usepackage{amssymb}
\usepackage{float}
\newtheorem{definition}{Definition}
\newtheorem{lemma}{Lemma}
\usepackage{cleveref}
\crefname{definition}{Definition}{Definitions}
\crefname{lemma}{Lemma}{Lemmas}
\title{Weak cosmic censorship for the circularly symmetric \\ Einstein-scalar field system in $2+1$ dimensions}
\author{Serban Cicortas}
\date{Source: arXiv:2605.19143v2 (CC BY 4.0)}
\begin{document}
\maketitle

\noindent\emph{Source and licence.} Weak cosmic censorship for the circularly symmetric \\ Einstein-scalar field system in $2+1$ dimensions. Version arXiv:2605.19143v2 (CC BY 4.0), distributed by arXiv under the Creative Commons Attribution 4.0 International licence, http://creativecommons.org/licenses/by/4.0/, as recorded in the arXiv licence metadata for this exact version and shown on the abstract page https://arxiv.org/abs/2605.19143v2. Full licence notice: \texttt{licenses/arxiv-2605.19143-CC-BY-4.0.txt}.\par\smallskip

\noindent\textit{Editorial scope.} Counted unit: the article's lemma 1 (source0.tex lines 1572--1583). The article's complete original proof of it is delivered with it (source0.tex lines 1584--1673, one \texttt{proof} environment of 90 lines). No mathematics is written, repaired or re-derived: the statement and the proof are the article's own lines, in the article's own notation, and no retained line is substituted or re-flowed.  Retained uncounted as the source-local context the proof needs: lines 527--533; lines 555--557; lines 1353--1355; lines 1359--1368; lines 1409--1411; lines 1457--1459; lines 1494--1496; lines 1513--1515; lines 1519--1521; lines 1565--1567.  Nothing was omitted from any retained span, so the package carries no omission record.  No retained line contains a citation, so the wrapper carries no bibliography and no bibliography entry.  No retained line contains a figure environment, an image-inclusion command or drawing code, so no image is delivered and the package declares no asset dependency.  Editorial formatting only, in addition: an AMS \texttt{amsart} preamble that restates the article's own declarations and macros used by the retained lines, namely the environments \texttt{definition}, \texttt{lemma} on separate counters, together with the standard packages those lines need.  The retained environments print in this document as Definition 1, Lemma 1 in order of appearance, whereas the article prints the same items with the numbering of its own full document.  The complete native source of the article as distributed in the arXiv e-print for this exact version is bundled unchanged as \texttt{sources/200896/source0.tex}.
\begin{align}
    2r\partial_u\partial_v\phi&=-\partial_ur\partial_v\phi-\partial_vr\partial_u\phi,\label{wave equation phi}\\
    2\partial_u\partial_vr&=\Lambda\Omega^2r,\label{wave equation r}\\
    2\partial_u\partial_v\log\Omega^{2}&=-2\partial_u\phi\partial_v\phi+\Lambda\Omega^2,\label{wave equation Omega}\\
    \partial_u\big(\Omega^{-2}\partial_ur\big)&=-\Omega^{-2}r(\partial_u\phi)^2,\label{Ray u}\\
    \partial_v\big(\Omega^{-2}\partial_vr\big)&=-\Omega^{-2}r(\partial_v\phi)^2\label{Ray v}.
\end{align}

\begin{equation}\label{dv log gamma}
    \partial_v\log\gamma=\frac{r(\partial_v\phi)^2}{\partial_vr}=\theta^2\partial_vr.
\end{equation}

\begin{equation}\label{gamma lower bound on C0-}
    \gamma(u,0)=\frac{-4\partial_ur(u,0)}{|\Lambda|r^2(u,0)+1-m(u,0)}\geq8|\Lambda|^{-1}|u|^{-2}.
\end{equation}

\begin{definition}\label{blueshift definition}
    We define the blueshift as the function $\mathbf{b}:[u_0,0)\times[0,v_1]\rightarrow[0,\infty):$
    \begin{equation}\label{blueshift definition eq}
        \mathbf{b}(u,v)=\int_{u_0}^u\frac{|\Lambda|}{2}r\gamma(u',v)du'.
    \end{equation}
    We denote $\mathbf{b}(u)=\mathbf{b}(u,0).$ We also define the function $\mathbf{s}:[u_0,0)\rightarrow[0,\infty):$
    \begin{equation}\label{blueshift strength definition}
        \mathbf{s}(u)=\int_{u_0}^u\frac{|\Lambda|}{4}r^2\gamma(u',0)du'.
    \end{equation}
\end{definition}

    \begin{equation}\label{background energy bound 1}
    \int_0^v\theta^2\partial_vr(u,v')dv'\leq10^4|\Lambda|\cdot v|u|e^{-\mathbf{b}(u)}\gamma(u,0),
    \end{equation}

\begin{equation}\label{dv r in terms of b(u)}
    \partial_vr(u,v)=\frac{1}{2}e^{-\mathbf{b}(u,v)}=\frac{1}{2}e^{-\mathbf{b}(u)}\big(1+O(\epsilon)\big)\leq 10e^{-\mathbf{b}(u)},
\end{equation}

    \begin{equation}\label{definition of P lambda}
    \mathcal{P}_{\lambda}=\bigg\{u\in[u_0,0),\ 0\leq v\leq\min\big(v_0',\lambda^4,|u|^{\frac{1}{n}}\big),\ |\lambda| v^{n+\frac{1}{4}}\leq e^{-\mathbf{b}(u)}|u|^{\frac{1}{2}+\frac{1}{10n}}\sqrt{\gamma(u,0)},\ v^{\frac{1}{10n}}\mathbf{s}(u)\leq\epsilon'\bigg\},
    \end{equation}

\begin{equation}\label{A1 blueshift instability}
    \partial_vr_{\lambda}(u,v)\leq 10e^{-\mathbf{b}(u)},
\end{equation}

\begin{equation}\label{A3 blueshift instability}
    \int_0^v\big|\theta_{\lambda}-\theta\big|^2(u,v')dv'\leq 100\lambda^2v^{2n+1}e^{2\mathbf{b}(u)}.
\end{equation}

\begin{equation}\label{dv r lambda in terms of b(u)}
    \partial_vr_{\lambda}(u,v)=\frac{1}{2}e^{-\mathbf{b}_{\lambda}(u,v)}=\frac{1}{2}e^{-\mathbf{b}(u)}\big(1+O(v^{\frac{1}{4}})\big),
\end{equation}

\begin{lemma}
    Under the bootstrap assumptions \eqref{A1 blueshift instability}-\eqref{A3 blueshift instability}, the following bounds hold for all $(u,v)\in[u_0,\Tilde{u}]\times[0,\Tilde{v}]$:
    \begin{align}
        \big|\widetilde{\gamma}\big|(u,v)\lesssim& ve^{-\mathbf{b}(u)}|u|\gamma^2(u,0)\big[v\mathbf{s}(u)\big]^{n}+|\lambda| v^{n+1}\gamma(u,0)\sqrt{|u|\gamma(u,0)}+\lambda^2v^{2n+1}\gamma(u,0)e^{\mathbf{b}(u)},\label{blueshift lemma bound 1}\\
        \big|\widetilde{\partial_vr}\big|(u,v)\lesssim& e^{-\mathbf{b}(u)}\big[v\mathbf{s}(u)\big]^{n+1}+|\lambda|v^{n+1}e^{-\mathbf{b}(u)}\mathbf{x}(u)+\lambda^2v^{2n+1},\label{blueshift lemma bound 2}\\
        \big|\widetilde{r}\big|(u,v)\lesssim & e^{-\mathbf{b}(u)}v\big[v\mathbf{s}(u)\big]^{n+1}+|\lambda| v^{n+2}e^{-\mathbf{b}(u)}\mathbf{x}(u)+\lambda^2v^{2n+2},\label{blueshift lemma bound 3}\\
        \big|\widetilde{\partial_ur}\big|(u,v)\lesssim& v|u|\big[v\mathbf{s}(u)\big]^{n}+|\lambda| v^{n+1}e^{-\mathbf{b}(u)/2}|u|^{1/2}\gamma(u,0)+\lambda^2 v^{2n+1}\gamma(u,0),\label{blueshift lemma bound 4}\\
        \big|\widetilde{\sqrt{r}\partial_u\phi}\big|(u,v)\lesssim& v|u|^{-1/2}e^{-\mathbf{b}(u)}\sqrt{\gamma(u,0)}\big[v\mathbf{s}(u)\big]^{n}+|\lambda|\cdot|u|^{-1}v^{n+1},\label{blueshift lemma bound 5}
    \end{align}
    where we use the notation:
    \[\mathbf{x}(u)=\int_{u_0}^u\big(|u'|\gamma(u',0)\big)^{3/2}du'.\]
\end{lemma}

\begin{proof}
    We notice that we have the inequality:
    \[\big|\widetilde{\theta^2\partial_vr}\big|\lesssim\theta^2\big|\widetilde{\partial_vr}\big|+e^{-\mathbf{b}(u)}|\theta|\cdot\big|\widetilde{\theta}\big|+e^{-\mathbf{b}(u)}\big|\widetilde{\theta}\big|^2.\]
    We integrate this inequality, and use \eqref{background energy bound 1}, \eqref{dv r in terms of b(u)}, \eqref{dv r lambda in terms of b(u)} for the first term, \eqref{background energy bound 1}, \eqref{A3 blueshift instability} for the second term, and \eqref{A3 blueshift instability} for the third term. Thus, we get that for all $(u,v)\in[u_0,\Tilde{u}]\times[0,\Tilde{v}]$:
    \begin{align}\label{preliminary bound difference of theta}
        \int_0^v\big|\widetilde{\theta^2\partial_vr}\big|(u,v')dv'&\lesssim v|u|\gamma(u,0)\sup_{v'\in[0,v]}\big|\widetilde{\partial_vr}\big|(u,v')+|\lambda| v^{n+1}\sqrt{|u|\gamma(u,0)}+\lambda^2v^{2n+1}e^{\mathbf{b}(u)}\\
        &\lesssim v|u|\gamma(u,0)e^{-\mathbf{b}(u)}+|\lambda|v^{n+1}\sqrt{|u|\gamma(u,0)}+\lambda^2v^{2n+1}e^{\mathbf{b}(u)}.\notag
    \end{align}
    According to \eqref{definition of P lambda}, the RHS of \eqref{preliminary bound difference of theta} is bounded up to a constant by $v^{\frac{1}{2}}|u_0|^{1+\frac{1}{5n}}\Omega^2(u_0,0),$ so it is small for $v_0'$ sufficiently small in terms of $|\Lambda|,|u_0|,\Omega^2(u_0,0).$ As a result, \eqref{dv log gamma} implies that:
    \begin{align}\label{preliminary bound difference of gamma}
        \big|\widetilde{\gamma}\big|(u,v)&\lesssim \gamma(u,v)\cdot\int_0^v\big|\widetilde{\theta^2\partial_vr}\big|(u,v')dv'\\
        &\lesssim v|u|\gamma^2(u,0)e^{-\mathbf{b}(u)}+|\lambda|v^{n+1}\gamma(u,0)\sqrt{|u|\gamma(u,0)}+\lambda^2v^{2n+1}\gamma(u,0)e^{\mathbf{b}(u)}.\notag
    \end{align}
    
    By induction we prove that for any $0\leq i\leq n$ and $(u,v)\in[u_0,\Tilde{u}]\times[0,\Tilde{v}]$:
    \begin{align}
        \big|\widetilde{\gamma}\big|(u,v)\lesssim & ve^{-\mathbf{b}(u)}|u|\gamma^2(u,0)\big[v\mathbf{s}(u)\big]^{i}+|\lambda|v^{n+1}\gamma(u,0)\sqrt{|u|\gamma(u,0)}+\lambda^2v^{2n+1}\gamma(u,0)e^{\mathbf{b}(u)},\label{induction g}\\
        \big|\widetilde{\partial_vr}\big|(u,v)\lesssim & e^{-\mathbf{b}(u)}\big[v\mathbf{s}(u)\big]^{i}+|\lambda|v^{n+1}e^{-\mathbf{b}(u)}\mathbf{x}(u)+\lambda^2v^{2n+1}.\label{induction r}
    \end{align}
    According to \eqref{dv r in terms of b(u)}, \eqref{dv r lambda in terms of b(u)}, and \eqref{preliminary bound difference of gamma}, we already established the case $i=0$. We assume the induction hypothesis holds for $0\leq i\leq n-1$ and prove it for $i+1.$ We first have the preliminary estimate:
    \begin{equation}\label{preliminary bound difference of r}
        \big|\widetilde{r}\big|(u,v)\lesssim e^{-\mathbf{b}(u)}v\big[v\mathbf{s}(u)\big]^{i}+|\lambda| v^{n+2}e^{-\mathbf{b}(u)}\mathbf{x}(u)+\lambda^2v^{2n+2}.
    \end{equation}
    As a result, we get by combining \eqref{induction g} with \eqref{preliminary bound difference of r} and using \eqref{gamma lower bound on C0-} that for any $(u,v)\in[u_0,\Tilde{u}]\times[0,\Tilde{v}]$:
    \begin{align}\label{bound for tilde rgamma}
        \big|\widetilde{r\gamma}\big|(u,v)\lesssim & ve^{-\mathbf{b}(u)}|u|^2\gamma^2(u,0)\big[v\mathbf{s}(u)\big]^{i}+|\lambda| v^{n+1}\big(|u|\gamma(u,0)\big)^{3/2}+\lambda^2v^{2n+1}|u|\gamma(u,0)e^{\mathbf{b}(u)}\\
        & +|\lambda|\Omega^2(u_0,0) v^{n+2}\mathbf{x}(u)+\lambda^2\Omega^2(u_0,0)v^{2n+2}e^{\mathbf{b}(u)}.\notag
    \end{align}
    We notice that by \Cref{blueshift definition}, we have that:
    \begin{equation}\label{bound for the old y}
        \int_{u_0}^u|u'|\gamma(u',0)e^{\mathbf{b}(u')}du'\lesssim e^{\mathbf{b}(u)}.
    \end{equation}
    We integrate \eqref{bound for tilde rgamma}, using \eqref{bound for the old y} and the fact that $\mathbf{s},\mathbf{x}$ are non-decreasing to get that:
    \begin{equation}\label{preliminary bound difference of r gamma}
        \int_{u_0}^u\big|\widetilde{r\gamma}\big|(u',v)du'\lesssim \big[v\mathbf{s}(u)\big]^{i+1}+|\lambda| v^{n+1}\mathbf{x}(u)+\lambda^2v^{2n+1}e^{\mathbf{b}(u)}.
    \end{equation}
    We bound the last two terms on the RHS of \eqref{preliminary bound difference of r gamma} using \eqref{definition of P lambda}:
    \begin{align}
        |\lambda| v^{n+1}\mathbf{x}(u)&\lesssim\int_{u_0}^u|u'|^{\frac{3}{2}}\gamma(u',0)\cdot v^{\frac{3}{4}}|u'|^{\frac{1}{2}+\frac{1}{10n}}e^{-\mathbf{b}(u')}\gamma(u',0)du'\lesssim |u_0|^{\frac{1}{10n}}v^{\frac{3}{4}}\mathbf{s}(u),\label{preliminary bound x}\\
        \lambda^2v^{2n+1}e^{\mathbf{b}(u)}&\lesssim v^{\frac{1}{2}}e^{-\mathbf{b}(u)}|u|^{1+\frac{1}{5n}}\gamma(u,0)\lesssim|u_0|^{1+\frac{1}{5n}}v^{\frac{1}{2}}.\label{preliminary bound y}
    \end{align}
    In particular, \eqref{preliminary bound x}-\eqref{preliminary bound y} imply that the RHS of \eqref{preliminary bound difference of r gamma} is small, so using \eqref{wave equation r} we get:
    \begin{equation}\label{preliminary bound difference of dv r}
        \big|\widetilde{\partial_vr}\big|(u,v)\lesssim\partial_vr(u,v)\cdot\bigg|\exp\bigg(\int_{u_0}^u\frac{\Lambda}{2}\widetilde{r\gamma}(u',v)du'\bigg)-1\bigg|\lesssim e^{-\mathbf{b}(u)}\cdot\int_{u_0}^u\big|\widetilde{r\gamma}\big|(u',v)du'.
    \end{equation}
    Together with \eqref{preliminary bound difference of r gamma}, this implies \eqref{induction r} for $i+1.$ Similarly to the derivation of \eqref{preliminary bound difference of theta}, we have that:
    \begin{align}
        \int_0^v\big|\widetilde{\theta^2\partial_vr}\big|(u,v')dv'\lesssim&ve^{-\mathbf{b}(u)}|u|\gamma(u,0)\big[v\mathbf{s}(u)\big]^{i+1}+|\lambda|\cdot|u|\gamma(u,0)v^{n+2}e^{-\mathbf{b}(u)}\mathbf{x}(u)\label{mid preliminary bound difference of theta}\\
        &+\lambda^2|u|\gamma(u,0)v^{2n+2}+|\lambda| v^{n+1}\sqrt{|u|\gamma(u,0)}+\lambda^2v^{2n+1}e^{\mathbf{b}(u)},\notag
    \end{align}
    where the first three terms follow by using \eqref{induction r} for $i+1,$ and the last two terms are the same as in \eqref{preliminary bound difference of theta}. We bound the second term on the RHS of \eqref{mid preliminary bound difference of theta} using \eqref{definition of P lambda} as follows:
    \begin{align}\label{some bound I need later}
        |\lambda|\cdot|u|\gamma(u,0)v^{n+2}e^{-\mathbf{b}(u)}\mathbf{x}(u)&\lesssim|\lambda| v^{n+1}\sqrt{|u|\gamma(u,0)}\cdot v\sqrt{e^{-\mathbf{b}(u)}\gamma(u,0)}\int_{u_0}^u|u'|^2\gamma(u',0)\sqrt{e^{-\mathbf{b}(u')}\gamma(u',0)}du'\notag\\
        &\lesssim|\lambda| v^{n+1}\sqrt{|u|\gamma(u,0)}\cdot v\mathbf{s}(u)\lesssim|\lambda| v^{n+1}\sqrt{|u|\gamma(u,0)}.
    \end{align}
    Additionally, the third term on the RHS of \eqref{mid preliminary bound difference of theta} is bounded by the fifth term. Therefore, we can further bound the terms on the RHS of \eqref{mid preliminary bound difference of theta} to get for any $(u,v)\in[u_0,\Tilde{u}]\times[0,\Tilde{v}]$:
    \begin{equation}\label{better preliminary bound difference of theta}
        \int_0^v\big|\widetilde{\theta^2\partial_vr}\big|(u,v')dv'\lesssim ve^{-\mathbf{b}(u)}|u|\gamma(u,0)\big[v\mathbf{s}(u)\big]^{i+1}+|\lambda| v^{n+1}\sqrt{|u|\gamma(u,0)}+\lambda^2v^{2n+1}e^{\mathbf{b}(u)}.
    \end{equation}
    Arguing as we did for \eqref{preliminary bound difference of theta} above, the RHS of \eqref{better preliminary bound difference of theta} is small. We can repeat the derivation of \eqref{preliminary bound difference of gamma} and use \eqref{better preliminary bound difference of theta} instead of \eqref{preliminary bound difference of theta}. This implies that \eqref{induction g} holds for $i+1,$ completing our induction argument. 

    In particular, the estimate \eqref{induction g} for $i=n$ implies that \eqref{blueshift lemma bound 1} holds. Iterating the argument used in proving \eqref{preliminary bound difference of r}-\eqref{preliminary bound difference of dv r}, while starting with \eqref{induction g}-\eqref{induction r} for $i=n$, we also get \eqref{blueshift lemma bound 2} and the desired bound for $\widetilde{r}$ in \eqref{blueshift lemma bound 3}. Next, we have by \eqref{wave equation r}, \eqref{definition of P lambda}, and \eqref{blueshift lemma bound 1}-\eqref{blueshift lemma bound 3} that for any $(u,v)\in[u_0,\Tilde{u}]\times[0,\Tilde{v}]$:
    \begin{align*}
        \big|\widetilde{\partial_ur}\big|(u,v)&\lesssim\int_{0}^v\big|\widetilde{r\gamma\partial_vr}\big|(u,v')dv'\lesssim\int_{0}^v\bigg[\gamma(u,0)e^{-\mathbf{b}(u)}\big|\widetilde{r}\big|+|u|e^{-\mathbf{b}(u)}\big|\widetilde{\gamma}\big|+|u|\gamma(u,0)\big|\widetilde{\partial_vr}\big|\bigg](u,v')dv'\\
        &\lesssim v|u|\big[v\mathbf{s}(u)\big]^{n}+|\lambda|v^{n+2}e^{-\mathbf{b}(u)}|u|\gamma(u,0)\Big[\mathbf{x}(u)+\sqrt{|u|\gamma(u,0)}\Big]+\lambda^2 v^{2n+2}|u|\gamma(u,0)\\
        &\lesssim v|u|\big[v\mathbf{s}(u)\big]^{n}+|\lambda|v^{n+1}e^{-\mathbf{b}(u)/2}|u|^{1/2}\gamma(u,0)+\lambda^2 v^{2n+1}\gamma(u,0),
    \end{align*}
    which completes the proof of \eqref{blueshift lemma bound 4}. 

    In order to prove \eqref{blueshift lemma bound 5}, we first derive the following equation as a consequence of \eqref{wave equation phi}:
    \[\partial_v\big(\sqrt{r}\partial_u\phi\big)=-\frac{1}{2\sqrt{r}}\partial_ur\partial_v\phi=-\frac{1}{2r}\partial_ur\partial_vr\cdot\theta.\]
    We use this to get an equation for $\partial_v\big(\widetilde{\sqrt{r}\partial_u\phi}\big)$. Thus, we start with the following bound:
    \begin{equation}\label{preliminary bound difference du phi}
        \big|\widetilde{\sqrt{r}\partial_u\phi}\big|(u,v)\lesssim\int_0^v\bigg(|u|^{-1}\big|\widetilde{\partial_vr}\cdot\theta\big|+e^{-\mathbf{b}(u)}|u|^{-2}\big|\widetilde{r}\cdot\theta\big|+e^{-\mathbf{b}(u)}|u|^{-1}\Big(\big|\widetilde{\partial_ur}\cdot\theta\big|+\big|\widetilde{\theta}\big|\Big)\bigg)(u,v')dv'.
    \end{equation}
    We bound the terms on the RHS of \eqref{preliminary bound difference du phi} one by one. For the first term we have using \eqref{background energy bound 1}, \eqref{definition of P lambda}, \eqref{blueshift lemma bound 2}, and \eqref{some bound I need later}:
    \begin{align*}
        \int_0^v|u|^{-1}\big|\widetilde{\partial_vr}\cdot\theta\big|(u,v')dv'\lesssim &v|u|^{-1/2}e^{-\mathbf{b}(u)}\sqrt{\gamma(u,0)}\big[v\mathbf{s}(u)\big]^{n}+|u|^{-1}|\lambda|\sqrt{|u|\gamma(u,0)}v^{n+2}e^{-\mathbf{b}(u)}\mathbf{x}(u)\\
        &+\lambda^2v^{2n+2}|u|^{-1/2}\sqrt{\gamma(u,0)}\\
        \lesssim &v|u|^{-1/2}e^{-\mathbf{b}(u)}\sqrt{\gamma(u,0)}\big[v\mathbf{s}(u)\big]^{n}+|u|^{-1}|\lambda|v^{n+1}.
    \end{align*}
    Additionally, the second term on the RHS of \eqref{preliminary bound difference du phi} is bounded by the same quantity. This follows since $e^{-\mathbf{b}(u)}|u|^{-2}\lesssim1$ and the RHS of \eqref{blueshift lemma bound 3} is bounded by the RHS of \eqref{blueshift lemma bound 2}. For the third term on the RHS of \eqref{preliminary bound difference du phi} we use \eqref{background energy bound 1} and \eqref{blueshift lemma bound 4}:
    \begin{equation*}
        \int_0^ve^{-\mathbf{b}(u)}|u|^{-1}\big|\widetilde{\partial_ur}\cdot\theta\big|(u,v')dv'\lesssim v^2e^{-\mathbf{b}(u)}\sqrt{|u|\gamma(u,0)}\big[v\mathbf{s}(u)\big]^{n}+|\lambda|v^{n+2}+\lambda^2v^{2n+2}|u|^{-1/2}\sqrt{\gamma(u,0)}.
    \end{equation*}
    Finally, for the fourth term on the RHS of \eqref{preliminary bound difference du phi} we use the bootstrap assumption \eqref{A3 blueshift instability}:
    \[\int_0^ve^{-\mathbf{b}(u)}|u|^{-1}\big|\widetilde{\theta}\big|(u,v')dv'\lesssim\sqrt{v}e^{-\mathbf{b}(u)}|u|^{-1}\bigg(\int_0^v\big|\widetilde{\theta}\big|^2(u,v')dv'\bigg)^{1/2}\lesssim|u|^{-1}|\lambda|v^{n+1}.\]
    Combining the estimates for the terms on the RHS of \eqref{preliminary bound difference du phi}, we completed the proof of \eqref{blueshift lemma bound 5}.
    %For the fourth term on the RHS of \eqref{preliminary bound difference du phi} we use \eqref{background energy bound 1} and \eqref{blueshift lemma bound 3}:\begin{align*}\int_0^ve^{-\mathbf{b}(u)}|u|^{-2}\big|\widetilde{r}\cdot\theta\big|(u,v')dv'\lesssim &v^2e^{-\mathbf{b}(u)}\sqrt{|u|\gamma(u,0)}\big[v\mathbf{s}(u)\big]^{n}+|\lambda|\sqrt{|u|\gamma(u,0)}v^{n+3}e^{-\mathbf{b}(u)}\mathbf{x}(u)\\ &+\lambda^2v^{2n+3}\sqrt{|u|\gamma(u,0)}\\ \lesssim &v^2e^{-\mathbf{b}(u)}\sqrt{|u|\gamma(u,0)}\big[v\mathbf{s}(u)\big]^{n}+\lambda|v^{n+2}.\end{align*}
\end{proof}

\end{document}
